% Adapted from Gui-Ping (Roy) Liang 
% 編譯 / Compilation: XeLaTeX → Biber → XeLaTeX → XeLaTeX
% 本機指令 / Local command: latexmk -xelatex main-template.tex
\documentclass[11pt,a4paper]{article}
\usepackage[a4paper,margin=2.5cm]{geometry}
\usepackage{fontspec}
\usepackage{xeCJK}
% 中文字體與程式碼字體 / Chinese text and code fonts
\IfFontExistsTF{Noto Serif CJK TC}{
  \setCJKmainfont{Noto Serif CJK TC}
  \setCJKmonofont{Noto Serif CJK TC}
}{
  \IfFileExists{/System/Library/Fonts/Supplemental/Songti.ttc}{
    \setCJKmainfont[Path=/System/Library/Fonts/Supplemental/,FontIndex=7,BoldFont=Songti.ttc,BoldFeatures={FontIndex=2}]{Songti.ttc}
    \setCJKmonofont[Path=/System/Library/Fonts/Supplemental/,FontIndex=7,BoldFont=Songti.ttc,BoldFeatures={FontIndex=2}]{Songti.ttc}
  }{
    \setCJKmainfont{bsmi00lp.ttf}
    \setCJKmonofont{bsmi00lp.ttf}
  }
}
\usepackage{amsmath,amssymb}
\usepackage{graphicx} % 插入圖片 / Insert images with \includegraphics
\graphicspath{{images/}}
\usepackage{booktabs,tabularx}
\usepackage{setspace}
\setstretch{1.3}
\usepackage{enumitem}
\usepackage{xcolor}
\definecolor{darkblue}{RGB}{30,30,130}
\definecolor{darkred}{RGB}{130,10,10}
\definecolor{codegreen}{rgb}{0,0.6,0}
\definecolor{stefan_green}{RGB}{0,100,0}
\newcommand{\david}[1]{{\color{codegreen}#1}}
\newcommand{\change}[1]{{\color{stefan_green}#1}}
\newcommand{\notice}[1]{{\color{darkred}#1}}
\usepackage[backend=biber,style=authoryear,sorting=nyt]{biblatex}
\addbibresource{reference.bib}
\usepackage{hyperref}
\hypersetup{colorlinks=true,linkcolor=darkblue,citecolor=darkblue,urlcolor=darkblue}
\renewcommand{\abstractname}{摘要 / Abstract}
\renewcommand{\contentsname}{目錄 / Contents}
\renewcommand{\figurename}{圖 / Figure}
\renewcommand{\tablename}{表 / Table}
\renewcommand{\refname}{參考文獻 / References}

\title{\textbf{NTLDEV 7192 與 NTLDEV 7191}\\課程專用報告（研究紀要）模板\\[0.8em]\textbf{NTLDEV 7192 and NTLDEV 7191}\\Course Report Template/Research Note}
\author{姓名 / Name\thanks{系所、學號與電子郵件。Department, student ID, and email.}}
\date{\today}

\begin{document}
% 標題頁與摘要分頁 / Separate the title and abstract pages.
\maketitle
\clearpage
\begin{abstract}

[Write your abstract here. You can then comment out the text below using the percent sign.]


\begingroup
\color{red}
\noindent\textbf{注意：Overleaf 專案名稱。}請將 Overleaf 專案名稱改為「課程代碼-姓名-學號」，方便我在專案列表中辨識課程、學生姓名與學號。主檔維持 \texttt{main-template.tex} 或直接覆蓋\texttt{main.tex}均可。

\noindent\textbf{Important: Overleaf project name.} 

Rename your Overleaf project to \texttt{CourseCode-YourName-StudentID}. This helps the instructor identify your course, name, and student ID in the project list. You can keep the main file named \texttt{main-template.tex} or rename it to \texttt{main.tex}.

\par
\endgroup
\medskip

\noindent\textbf{篇幅提醒：}摘要至多 1,000 字，建議控制在 500 字以內，以便日後修改並用於研討會投稿。投稿時請以各研討會的摘要格式與篇幅規定為準。

\noindent\textbf{Length guidance:} The abstract must not exceed 1,000 words; preferably keep it within 500 words to facilitate adaptation for conference submissions. Follow the specific abstract format and length requirements of the target conference.

\medskip
本版模示範比較立法報告常用的文字格式、章節、圖片、表格、公式、交叉引用與文獻引用。
請將示範文字替換成自己的研究問題、資料、方法與發現。
所有示範數字均為教學用途，不代表實證結果。

This template demonstrates text formatting, sections, figures, tables, equations, cross-references, and citations for a comparative legislative studies report. Replace the examples with your research questions, data, methods, and findings. All numerical examples are for teaching and do not represent empirical results.
\end{abstract}
\clearpage
% 目錄獨立一頁 / Start the contents on a separate page.
\tableofcontents
\clearpage

\section{緒論 / Introduction}
\label{sec:introduction}
\subsection{研究動機與問題 / Motivation and Questions}
在此說明研究背景、問題意識與研究貢獻。例如：不同議會制度下，政黨如何維持議員的投票一致性？

Describe the background, research question, and contribution. For example: How do parties maintain voting unity under different parliamentary systems?

% 空一行表示新的段落；不要用連續的 \\ 代替段落。
% Start a paragraph with a blank line; do not use repeated line breaks.
這是一個新段落。使用 \textbf{粗體}強調關鍵概念，使用 \emph{emphasis}標示英文重點。
特殊字元需要跳脫，例如百分比 25\%、底線 party\_id、\& 與 \#。
註腳可以補充資料來源或定義。\footnote{這是註腳範例，請替換成必要的補充說明。 This is a sample footnote; replace it with a useful note.}

Use a blank line to start a new paragraph. Use \textbf{bold} for key concepts and \emph{emphasis} for emphasis. Escape special characters, such as 25\%, party\_id, \&, and \#. Footnotes can provide sources or definitions.

\subsection{條列與連結 / Lists and Links}
\begin{itemize}
  \item 研究對象：國家、議院或政黨。\par Cases: countries, chambers, or parties.
  \item 分析單位：議員、單次表決或議員與表決的組合。\par Unit of analysis: legislators, divisions, or legislator--division observations.
\end{itemize}
\begin{enumerate}[label=\arabic*.]
  \item 提出研究問題與理論預期。\par State the question and theoretical expectations.
  \item 整理資料，定義變數並進行分析。\par Prepare the data, define variables, and conduct the analysis.
  \item 回應問題並說明研究限制。\par Answer the question and discuss limitations.
\end{enumerate}
網址可使用 \url{https://www.parliament.uk/}，也可使用
\href{https://www.parliament.uk/}{自訂連結文字}。

Use \texttt{itemize} for bullet points and \texttt{enumerate} for numbered steps. Use \texttt{\textbackslash url} for a visible URL, or \texttt{\textbackslash href} for custom link text.

\section{西敏制下的投票 / Voting in the Westminster System}
\label{sec:theory}
% 文獻 key 必須與 reference.bib 中的 key 完全相同。
% Citation keys must exactly match the keys in reference.bib.
% \textcite：作者融入句子；\parencite：括號引用。
% Use \textcite for narrative citations and \parencite for parenthetical citations.
\textcite{cox1987}的著作可作為理解英國議會制度發展的閱讀起點。
括號式引用的寫法如下\parencite{cox1987}。
一次引用多筆文獻可寫為\parencite{cox1987,carey2009}。

The work of \textcite{cox1987} offers a starting point for reading about the development of British parliamentary institutions. Parenthetical citations look like this \parencite{cox1987}; multiple sources can be cited together \parencite{cox1987,carey2009}.

指定頁碼的語法為： / Syntax for citing specific pages:
\begin{quote}
\texttt{\textbackslash{}parencite[23--25]\{cox1987\}}
\end{quote}
以上頁碼僅示範語法；請核對原文後填入真正支持論點的頁碼。

The command above demonstrates a citation with page numbers. The pages are illustrative; check the source and enter the pages that support your claim.

\begin{quote}
這裡是區塊引文的排版示範，並非任何文獻的原文。
實際使用時應填入經核對的引文，並附作者、年份與頁碼。

\par This is a block quotation formatting example, not a quotation from a source. Replace it with a verified quotation and cite the author, year, and page.
\end{quote}

\subsection{研究假設 / Research Hypothesis}
\textbf{H1：}請在此寫出可由資料檢驗的理論預期，並交代理論理由。

\textbf{H1:} State a testable expectation and explain its theoretical basis.

\section{表決、跑票與辯論資料 / Data on Divisions, Defection and Debates}
\label{sec:data}
\subsection{資料來源與變數 / Data Sources and Variables}
在此交代研究期間、資料來源、樣本選擇、缺漏值處理與變數定義。
下列圖片是隨版模附上的占位圖，請換成自己的圖表。

Describe the study period, sources, sample selection, missing data treatment, and variable definitions. Replace the placeholder below with your own figure.

% figure 是浮動環境；[htbp] 表示優先放在此處、頁頂、頁底或浮動頁。
% figure is a float; [htbp] allows placement here, at the top, bottom, or on a float page.
% 寬度使用 \linewidth 的比例，通常比固定公分更容易維持版面。
% Scale image width relative to \linewidth to adapt to the available space.
\begin{figure}[htbp]
  \centering
  \includegraphics[width=0.65\linewidth]{example-image.pdf}
  \caption{研究架構占位圖 / Research Framework Placeholder}
  \label{fig:framework} % label 放在 caption 後方 / Put label after caption
  \par\small 資料來源：LaTeX mwe 套件示範圖，非研究資料。\\Source: Example image from the LaTeX mwe package; not research data.
\end{figure}
如圖~\ref{fig:framework}所示，可以在內文引用圖片編號。
換圖時將 PNG、JPG 或 PDF 放入 \texttt{images/}，再修改
\texttt{\textbackslash{}includegraphics} 的檔名。

Refer to Figure~\ref{fig:framework} in the text. To replace the image, put a PNG, JPG, or PDF in \texttt{images/} and update the filename in \texttt{\textbackslash includegraphics}.

\begin{table}[htbp]
  \centering
  \caption{變數定義與示範值 / Variable Definitions and Illustrative Values}
  \label{tab:variables}
  \begin{tabularx}{\linewidth}{>{\raggedright\arraybackslash}p{0.28\linewidth}Xr}
    \toprule
    變數 / Variable & 定義 / Definition & 值 / Value \\
    \midrule
    跑票次數\newline Defections & 議員投票與所屬政黨立場不一致的次數。\newline Votes against the party position. & 8 \\
    參與表決次數\newline Votes cast & 納入分析的有效表決次數。\newline Valid votes included in the analysis. & 100 \\
    跑票比例\newline Defection rate & 跑票次數除以參與表決次數。\newline Defections divided by votes cast. & 0.08 \\
    \bottomrule
  \end{tabularx}
  \par\small 註：所有數字僅用於排版練習。\\Note: All values are fictional and used for formatting practice.
\end{table}
表~\ref{tab:variables}示範三線表；欄位中的 \texttt{l}、\texttt{r} 與 \texttt{X}
分別表示靠左、靠右與可自動換行的欄位。

Table~\ref{tab:variables} demonstrates a three-rule table. The column types \texttt{l}, \texttt{r}, and \texttt{X} mean left-aligned, right-aligned, and automatically wrapping columns.

\section{跑票分析 / Analysis of Rebellion}
\label{sec:analysis}
\subsection{公式排版 / Formatting Equations}
行內公式以美元符號包住，例如 $r_i = d_i/n_i$。
需要編號時使用 \texttt{equation} 環境：

Enclose inline mathematics in dollar signs, for example $r_i=d_i/n_i$. Use the \texttt{equation} environment for a numbered equation:
\begin{equation}
  r_i = \frac{d_i}{n_i}, \qquad n_i > 0.
  \label{eq:rebellion}
\end{equation}
式~\eqref{eq:rebellion}中，$d_i$ 為議員 $i$ 的跑票次數，$n_i$ 為參與表決次數。
若沒有有效表決，應另行說明處理方式。

In Equation~\eqref{eq:rebellion}, $d_i$ is the number of defections by legislator $i$, and $n_i$ is the number of votes cast. Explain how observations with no valid votes are handled.

多行公式可使用 \texttt{align} 環境，以 \texttt{\&} 對齊：

Use \texttt{align} for multiline equations, with \texttt{\&} marking alignment points:
\begin{align}
  \eta_i &= \alpha + \beta_1 x_{1i} + \beta_2 x_{2i}, \\
  p_i &= \frac{\exp(\eta_i)}{1+\exp(\eta_i)}.
\end{align}
以上僅為公式排版示範；實際模型應由研究問題與資料結構決定。

These equations demonstrate formatting only; choose an actual model based on the research question and data structure.

\subsection{結果與交叉引用 / Results and Cross-References}
可引用第~\ref{sec:theory}~節的理論、第~\ref{sec:data}~節的資料，
或指出圖~\ref{fig:framework}位於第~\pageref{fig:framework}~頁。
請勿手動輸入圖表與公式編號，以免增刪內容後編號失去一致性。

Refer to the theory in Section~\ref{sec:theory}, the data in Section~\ref{sec:data}, or Figure~\ref{fig:framework} on page~\pageref{fig:framework}. Use cross-reference commands so numbering stays correct when content changes.

\section{結論 / Conclusion}
在此回應研究問題，概述主要發現、研究限制與未來研究方向。
撰寫正式報告時，請移除本版模的操作說明與練習內容。

Answer the research question, summarize the findings, and discuss limitations and future research. Remove the instructional text and exercises when preparing the final report.



\clearpage

\printbibliography[heading=bibintoc,title={參考文獻 / References}]


\appendix
\section{常用指令速查表 / Common Command Quick Reference}
依用途查找下列指令，將範例中的文字、檔名與標籤替換為自己的內容。

Find the command you need below and replace the sample text, filenames, and labels with your own.

\begingroup
\small
\setlength{\tabcolsep}{4pt}
\renewcommand{\arraystretch}{1.35}
\noindent\begin{tabularx}{\linewidth}{@{}>{\raggedright\arraybackslash}p{0.30\linewidth}X@{}}
\toprule
用途 / Purpose & 指令 / Command \\
\midrule
章節 / Section & \texttt{\textbackslash{}section\{Title\}} \\
小節 / Subsection & \texttt{\textbackslash{}subsection\{Title\}} \\
次小節 / Subsubsection & \texttt{\textbackslash{}subsubsection\{Title\}} \\
粗體 / Bold & \texttt{\textbackslash{}textbf\{Text\}} \\
強調 / Emphasis & \texttt{\textbackslash{}emph\{Text\}} \\
新段落 / New paragraph & 原始碼空一行。 / Leave a blank line in the source. \\
註腳 / Footnote & \texttt{\textbackslash{}footnote\{Note\}} \\
網址 / URL & \texttt{\textbackslash{}url\{https://example.org\}} \\
連結文字 / Link text & \texttt{\textbackslash{}href\{https://example.org\}\{Text\}} \\
特殊字元 / Special characters & \texttt{\textbackslash\%\ \textbackslash\&\ \textbackslash\_\ \textbackslash\#\ \textbackslash\$\ \textbackslash\{\ \textbackslash\}} \\
\bottomrule
\end{tabularx}

\medskip
\noindent\textbf{圖片、表格與交叉引用 / Figures, Tables, and Cross-References}

\noindent\begin{tabularx}{\linewidth}{@{}>{\raggedright\arraybackslash}p{0.30\linewidth}X@{}}
\toprule
用途 / Purpose & 指令 / Command \\
\midrule
插入圖片 / Insert image & \texttt{\textbackslash{}includegraphics[width=.7\textbackslash{}linewidth]\{file.pdf\}} \\
圖表標題 / Caption & \texttt{\textbackslash{}caption\{Caption\}} \\
設定標籤 / Set label & \texttt{\textbackslash{}label\{fig:example\}} \\
引用圖表／章節 / Reference & \texttt{\textbackslash{}ref\{fig:example\}} \\
引用所在頁 / Page reference & \texttt{\textbackslash{}pageref\{fig:example\}} \\
表格水平線 / Table rules & \texttt{\textbackslash{}toprule} / \texttt{\textbackslash{}midrule} / \texttt{\textbackslash{}bottomrule} \\
分隔欄／結束列 / Column / row & \texttt{\&} / \texttt{\textbackslash{}\textbackslash{}} \\
\bottomrule
\end{tabularx}

\medskip
\noindent\begin{minipage}{\linewidth}
\textbf{文獻與公式 / Citations and Equations}\par\smallskip

\noindent\begin{tabularx}{\linewidth}{@{}>{\raggedright\arraybackslash}p{0.30\linewidth}X@{}}
\toprule
用途 / Purpose & 指令 / Command \\
\midrule
作者式引用 / Narrative citation & \texttt{\textbackslash{}textcite\{cox1987\}} \\
括號式引用 / Parenthetical citation & \texttt{\textbackslash{}parencite\{cox1987\}} \\
附頁碼引用 / Cite pages & \texttt{\textbackslash{}parencite[23--25]\{cox1987\}} \\
多篇引用 / Multiple sources & \texttt{\textbackslash{}parencite\{cox1987,carey2009\}} \\
行內公式 / Inline math & \texttt{\$r\_i = d\_i/n\_i\$} \\
分數／上下標 / Fraction / scripts & \texttt{\textbackslash{}frac\{a\}\{b\}} / \texttt{x\_i} / \texttt{x\textasciicircum{}2} \\
公式引用 / Equation reference & \texttt{\textbackslash{}eqref\{eq:example\}} \\
\bottomrule
\end{tabularx}
\end{minipage}
\endgroup

\subsection*{常用環境 / Common Environments}
下列區塊可直接複製。 / Copy and adapt the blocks below.

\noindent\textbf{條列與編號 / Bulleted and Numbered Lists}
\begin{verbatim}
\begin{itemize}          % 條列 / Bullets
  \item First point
  \item Second point
\end{itemize}

\begin{enumerate}        % 編號 / Numbers
  \item First step
  \item Second step
\end{enumerate}
\end{verbatim}

\noindent\textbf{圖片 / Figure}
\begin{verbatim}
\begin{figure}[htbp]
  \centering
  \includegraphics[width=.7\linewidth]{example-image.pdf}
  \caption{圖說 / Figure caption}
  \label{fig:example}
\end{figure}
\end{verbatim}

\noindent\begin{minipage}{\linewidth}
\textbf{表格 / Table}
\begin{verbatim}
\begin{table}[htbp]
  \centering
  \caption{表格標題 / Table caption}
  \label{tab:example}
  \begin{tabular}{lr}
    \toprule
    Item & Value \\
    \midrule
    A & 10 \\
    B & 20 \\
    \bottomrule
  \end{tabular}
\end{table}
\end{verbatim}
\end{minipage}\par\medskip

\noindent\textbf{編號公式 / Numbered Equation}
\begin{verbatim}
\begin{equation}
  r_i = \frac{d_i}{n_i}
  \label{eq:example}
\end{equation}
\end{verbatim}

\noindent\textbf{新增文獻 / Adding a Reference}
將下列格式加入 \texttt{reference.bib}，填入真實文獻資訊，再於內文使用相同的 key 引用。

Add an entry to \texttt{reference.bib}, replace the placeholder metadata with the actual source details, and cite the same key in the text.
\begin{verbatim}
@book{author2026,
  author    = {Surname, Given Name},
  title     = {Book Title},
  year      = {2026},
  publisher = {Publisher}
}
\end{verbatim}

\noindent\textbf{提醒 / Reminders}
\begin{itemize}
  \item 圖表的標籤放在圖說之後；文獻 key 須與資料庫一致。\par
  Put figure and table labels after captions; citation keys must match the bibliography.
  \item 範例頁碼須核對原文；圖片放在 \texttt{images/}。\par
  Verify citation pages against the source; store images in \texttt{images/}.
  \item 新增引用後重新編譯。\par
  Recompile after adding references.
\end{itemize}
% 預設只列出已引用的文獻；若要列出 bib 中全部文獻，可啟用下一行。
% Only cited works are listed by default; uncomment the next line to list all entries.
% \nocite{*}
\end{document}
